\section{Lexicon Reference}\label{app:lexicon}

\subsection{Composition}

The lexicon contains \LexiconTotal\ entries in two tiers, and the distinction is
load-bearing:

\begin{description}[leftmargin=1.4em, style=nextline, font=\ttfamily]
  \item[corpus (\LexiconCorpus\ entries)] Each one cites the statement identifiers in which
    it was observed. \texttt{LexiconSpec} validation \emph{rejects} a corpus entry with an
    empty evidence list, so this tier cannot be padded.
  \item[dictionary (\LexiconDictionary\ entries)] Closed-class French and Francanglais
    vocabulary sourced from reference material --- determiners, pronouns, prepositions,
    conjunctions and common verb forms. These carry no corpus evidence and make no claim of
    attestation.
\end{description}

\begin{limitbox}[A claim we will not make]
Earlier design documentation for this project stated that every lexicon entry existed only
because it was attested in a reviewed Yaound\'e statement. That is not true of the current
lexicon and it is corrected here: the majority of entries are dictionary-derived. The two
tiers are kept separately labelled precisely so that the distinction survives into any
report generated from the data.
\end{limitbox}

\subsection{Corpus-attested entries}

The table below lists the corpus tier in full, with the statement identifiers that justify
each entry. The dictionary tier is omitted for length and is browsable in the application's
grammar view, which supports incremental search.

\begin{longtable}{@{}llll p{0.30\linewidth}@{}}
\toprule
\textbf{Canonical} & \textbf{Terminal} & \textbf{POS} & \textbf{Origin} & \textbf{Attested in} \\
\midrule
\endfirsthead
\toprule
\textbf{Canonical} & \textbf{Terminal} & \textbf{POS} & \textbf{Origin} & \textbf{Attested in} \\
\midrule
\endhead
\bottomrule
\endfoot
\texttt{combi} & NOUN & noun & uncertain & \small demo-001, demo-002, demo-007, \textbackslash{}ldots \\
\texttt{on} & PRON & pronoun & french & \small demo-001, demo-007, demo-009, \textbackslash{}ldots \\
\texttt{go} & VERB & verb & english & \small demo-001, demo-010 \\
\texttt{au} & PREP & preposition & french & \small demo-001, demo-002, demo-003, \textbackslash{}ldots \\
\texttt{kwatt} & NOUN & noun & uncertain & \small demo-001, demo-003, demo-008, \textbackslash{}ldots \\
\texttt{en} & PREP & preposition & french & \small demo-001 \\
\texttt{taco} & NOUN & noun & uncertain & \small demo-001 \\
\texttt{j'ai} & VERB & verb & french & \small demo-001 \\
\texttt{les} & DET & determiner & french & \small demo-001, demo-011 \\
\texttt{dos} & NOUN & noun & french & \small demo-001, demo-011 \\
\texttt{pour} & PREP & preposition & french & \small demo-001, demo-004 \\
\texttt{la} & DET & determiner & french & \small demo-001, demo-005, demo-008, \textbackslash{}ldots \\
\texttt{course} & NOUN & noun & french & \small demo-001 \\
\texttt{le} & DET & determiner & french & \small demo-002, demo-003, demo-004, \textbackslash{}ldots \\
\texttt{réseau} & NOUN & noun & french & \small demo-002, demo-012 \\
\texttt{ne} & NEG & negation\_particle & french & \small demo-002, demo-008, demo-012 \\
\texttt{passe} & VERB & verb & french & \small demo-002, demo-012 \\
\texttt{pas} & NEG & negation\_particle & french & \small demo-002, demo-004, demo-008, \textbackslash{}ldots \\
\texttt{je} & PRON & pronoun & french & \small demo-002, demo-004, demo-005, \textbackslash{}ldots \\
\texttt{peux} & VERB & verb & french & \small demo-002 \\
\texttt{même} & ADV & adverb & french & \small demo-002 \\
\texttt{send} & VERB & verb & english & \small demo-002 \\
\texttt{message} & NOUN & noun & french & \small demo-002 \\
\texttt{ils} & PRON & pronoun & french & \small demo-003 \\
\texttt{ont} & VERB & verb & french & \small demo-003 \\
\texttt{encore} & ADV & adverb & french & \small demo-003, demo-011 \\
\texttt{coupé} & VERB & verb & french & \small demo-003 \\
\texttt{courant} & NOUN & noun & french & \small demo-003, demo-012 \\
\texttt{wanda} & VERB & verb & uncertain & \small demo-003 \\
\texttt{baisse} & VERB & verb & french & \small demo-004 \\
\texttt{un} & DET & determiner & french & \small demo-004 \\
\texttt{peu} & NOUN & noun & french & \small demo-004 \\
\texttt{prix} & NOUN & noun & french & \small demo-004, demo-007 \\
\texttt{des} & DET & determiner & french & \small demo-004, demo-009 \\
\texttt{tomates} & NOUN & noun & french & \small demo-004 \\
\texttt{n'ai} & VERB & verb & french & \small demo-004 \\
\texttt{assez} & ADV & adverb & french & \small demo-004 \\
\texttt{de} & PREP & preposition & french & \small demo-004, demo-005, demo-010 \\
\texttt{nkap} & NOUN & noun & uncertain & \small demo-004 \\
\texttt{buy} & VERB & verb & english & \small demo-004 \\
\texttt{tout} & ADV & adverb & french & \small demo-004 \\
\texttt{ça} & PRON & pronoun & french & \small demo-004 \\
\texttt{avec} & PREP & preposition & french & \small demo-005 \\
\texttt{cette} & DET & determiner & french & \small demo-005 \\
\texttt{pluie} & NOUN & noun & french & \small demo-005 \\
\texttt{préfère} & VERB & verb & french & \small demo-005 \\
\texttt{nang} & VERB & verb & uncertain & \small demo-005 \\
\texttt{à} & PREP & preposition & french & \small demo-005, demo-006, demo-012 \\
\texttt{piol} & NOUN & noun & uncertain & \small demo-005, demo-008, demo-012 \\
\texttt{waka} & VERB & verb & uncertain & \small demo-005, demo-008 \\
\texttt{dehors} & ADV & adverb & french & \small demo-005 \\
\texttt{taximan} & NOUN & noun & mixed & \small demo-006, demo-011 \\
\texttt{est} & VERB & verb & french & \small demo-006, demo-012 \\
\texttt{parti} & VERB & verb & french & \small demo-006 \\
\texttt{fala} & VERB & verb & uncertain & \small demo-006 \\
\texttt{l'essence} & NOUN & noun & french & \small demo-006, demo-011 \\
\texttt{sa} & DET & determiner & french & \small demo-006 \\
\texttt{caisse} & NOUN & noun & french & \small demo-006 \\
\texttt{viens} & VERB & verb & french & \small demo-007 \\
\texttt{tchop} & VERB & verb & uncertain & \small demo-007 \\
\texttt{poisson} & NOUN & noun & french & \small demo-007 \\
\texttt{ici} & ADV & adverb & french & \small demo-007 \\
\texttt{peut} & VERB & verb & french & \small demo-007 \\
\texttt{discuter} & VERB & verb & french & \small demo-007 \\
\texttt{bendskin} & NOUN & noun & mixed & \small demo-008 \\
\texttt{va} & VERB & verb & french & \small demo-008 \\
\texttt{me} & PRON & pronoun & french & \small demo-008 \\
\texttt{bring} & VERB & verb & english & \small demo-008 \\
\texttt{veux} & VERB & verb & french & \small demo-008 \\
\texttt{jusqu'à} & PREP & preposition & french & \small demo-008 \\
\texttt{m'a} & VERB & verb & french & \small demo-009 \\
\texttt{djoss} & VERB & verb & uncertain & \small demo-009, demo-010 \\
\texttt{qu'il} & CONJ & conjunction & french & \small demo-009 \\
\texttt{y} & ADV & adverb & french & \small demo-009 \\
\texttt{a} & VERB & verb & french & \small demo-009 \\
\texttt{vols} & NOUN & noun & french & \small demo-009 \\
\texttt{garde} & VERB & verb & french & \small demo-009 \\
\texttt{bien} & ADV & adverb & french & \small demo-009 \\
\texttt{ton} & DET & determiner & french & \small demo-009 \\
\texttt{phone} & NOUN & noun & english & \small demo-009 \\
\texttt{arrête} & VERB & verb & french & \small demo-010 \\
\texttt{school} & NOUN & noun & english & \small demo-010 \\
\texttt{prof} & NOUN & noun & french & \small demo-010 \\
\texttt{nous} & PRON & pronoun & french & \small demo-010 \\
\texttt{wait} & VERB & verb & english & \small demo-010 \\
\texttt{déjà} & ADV & adverb & french & \small demo-010 \\
\texttt{demande} & VERB & verb & french & \small demo-011 \\
\texttt{il} & PRON & pronoun & french & \small demo-011 \\
\texttt{dit} & VERB & verb & french & \small demo-011 \\
\texttt{que} & CONJ & conjunction & french & \small demo-011 \\
\texttt{coûte} & VERB & verb & french & \small demo-011 \\
\texttt{trop} & ADV & adverb & french & \small demo-011 \\
\texttt{cher} & ADJ & adjective & french & \small demo-011 \\
\texttt{back} & ADJ & adjective & english & \small demo-012 \\
\texttt{mais} & CONJ & conjunction & french & \small demo-012 \\
\texttt{toujours} & ADV & adverb & french & \small demo-012 \\
\texttt{à sec} & ADJ & adjective\_phrase & french & \small demo-006 \\
\texttt{au lieu de} & PREP & preposition\_phrase & french & \small demo-005 \\
\texttt{ashia} & INTJ & interjection & pidgin & \small demo-013 \\
\texttt{on dit quoi} & FORMULA & greeting\_formula & french & \small demo-014 \\
\texttt{suis} & VERB & verb & french & \small demo-016 \\
\texttt{fin} & NOUN & noun & french & \small demo-016 \\
\texttt{du} & PREP & preposition & french & \small demo-016 \\
\texttt{mois} & NOUN & noun & french & \small demo-016 \\
\texttt{ma} & DET & determiner & french & \small demo-019 \\
\texttt{pousse} & VERB & verb & french & \small demo-022 \\
\end{longtable}


\subsection{Duplicate prevention}

Canonical forms are unique across the whole lexicon, enforced at specification load time.
This caught two genuine collisions during development, where the same slang term was added
independently in two different vocabulary batches --- the kind of error that would otherwise
have produced a silently shadowed entry and an unexplainable parse result.
